\documentclass[11pt]{article}
\usepackage[margin=1in]{geometry}
\usepackage{hyperref}
\usepackage{enumitem}
\usepackage{booktabs}
\usepackage{amsmath}
\title{Observation Bandwidth, Reobservation, and Causal Order\\
\large A Falsifiable Experimental Program Derived from the QIK-VRT Computational Architecture}
\author{Ingolf Lohmann}
\date{August 2026}

\begin{document}
\maketitle

\begin{abstract}
QIK-VRT separates temporal sequence from causal order, execution from observed effect, change from improvement, and monitoring from complete observation. A bounded repository-native precursor executed adaptive sampling tests, a real Firefox-mediated prepare/commit roundtrip, post-effect backend reobservation, and exact-head causal/Motorola-68000 source tests. The precursor is computational evidence, not a physical result. This paper registers five falsifiable hypotheses concerning observation-rate adequacy, false completion reduction through reobservation, causal reconstruction under clock disorder, receive-order inversion without causal inversion, and incremental predictive value of causal-order features in two independently instrumented physical apparatuses. The stronger identity claim that physical time equals QIK-VRT causal time remains held because it lacks a discriminating operationalization.
\end{abstract}

\section{Scope}
The central invariant is
\[
\mathrm{CAUSALITY} \neq \mathrm{SEQUENCE}.
\]
The experimental program asks whether timestamp order, receive order, explicit predecessor order, intervention order, protocol acknowledgement, and post-effect reobservation can be empirically distinguished under controlled perturbations.

\section{Computational precursor}
The exact source head is \texttt{a7cea28de6ab435c01211d522613fb811bfd91b2}, tree \texttt{2ef860cfebd6005d7993818d2881b46b7dcf3212}, workflow run 32536665914, job 96938778530. The receipt digest is \texttt{f081469cdab77d951adb60a406d8012e5e773368fa92d850a2f499981838dacf}. Its scope is \texttt{BOUNDED\_LOOPBACK\_TERMINAL\_INPUT\_ONLY}; it records no external effect, no physical Mega ST execution, and no general \texttt{EFFECT\_ACK\_DONE}.

\section{Registered hypotheses}
\begin{description}[leftmargin=2.2cm,style=nextline]
\item[H-OBS-01] Above-bound observation reduces missed-transition and aliasing error for a registered finite-bandwidth transition generator.
\item[H-EFF-02] Independent post-effect reobservation reduces false completion classifications relative to acknowledgement alone under injected faults.
\item[H-CAU-03] Explicit predecessor bindings improve intervention-grounded causal reconstruction over timestamp order under skew, delay, and reordering.
\item[H-ORD-04] Observer-relative receive-order inversion can preserve one intervention-defined causal graph.
\item[H-PHY-05] Causal-order features add replicated out-of-sample predictive information beyond timestamps in two independently instrumented physical apparatuses.
\end{description}

The controlling nulls, endpoints, sample minima, decision rules, and falsification conditions are machine-readable in \texttt{HYPOTHESIS\_REGISTRY.json}.

\section{Blocked identity claim}
The statement
\[
\mathrm{PHYSICAL\_TIME} = \mathrm{QIKVRT\_CAUSAL\_TIME}
\]
is not an active hypothesis. It remains \texttt{HOLD\_NO\_OPERATIONALIZATION} until a measurement distinguishes it from weaker predictive, representational, or existing-physics alternatives.

\section{Protocol}
The protocol separates a synthetic observation benchmark, a fault-injected effect system, a distributed causal-order benchmark, and a two-apparatus physical bridge. Development data are separated from frozen held-out evaluation. Equal-information comparisons are mandatory. Negative results remain valid outcomes.

\section{Boundaries}
A successful engineering or causal-order experiment would support only its registered scope. Even a replicated physical-bridge result would not establish a new fundamental law, physical retrocausality, quantum causal order, machine consciousness, or general effect completion.

\section{Conclusion}
The transition from software architecture to science is:
\[
\mathrm{COMPUTATIONAL\ RESULT}
\rightarrow
\mathrm{OPERATIONAL\ DEFINITION}
\rightarrow
\mathrm{REGISTERED\ HYPOTHESIS}
\rightarrow
\mathrm{CONTROLLED\ INTERVENTION}
\rightarrow
\mathrm{INDEPENDENT\ OBSERVATION}
\rightarrow
\mathrm{FALSIFICATION\ OR\ SUPPORT}.
\]
Repository evidence motivates this chain but does not substitute for physical evidence.

\begin{thebibliography}{9}
\bibitem{nyquist1928}
H. Nyquist, ``Certain Topics in Telegraph Transmission Theory,'' \emph{Transactions of the AIEE}, vol. 47, no. 2, pp. 617--644, 1928. doi:10.1109/T-AIEE.1928.5055024.

\bibitem{shannon1949}
C. E. Shannon, ``Communication in the Presence of Noise,'' \emph{Proceedings of the IRE}, vol. 37, no. 1, pp. 10--21, 1949. doi:10.1109/JRPROC.1949.232969.

\bibitem{lamport1978}
L. Lamport, ``Time, Clocks, and the Ordering of Events in a Distributed System,'' \emph{Communications of the ACM}, vol. 21, no. 7, pp. 558--565, 1978. doi:10.1145/359545.359563.

\bibitem{bombelli1987}
L. Bombelli, J. Lee, D. Meyer, and R. D. Sorkin, ``Space-Time as a Causal Set,'' \emph{Physical Review Letters}, vol. 59, pp. 521--524, 1987. doi:10.1103/PhysRevLett.59.521.
\end{thebibliography}

\end{document}
